% Insert after the conclusion and before the bibliography.
% If the manuscript already starts its appendices elsewhere, omit the next line.
\appendices
\section{Real-Coordinate Fisher Information and Geometric CCRB}
\label{app:revision-ccrb}

\subsection{Channel coordinates and measurement information}

For receive element $n$, define
\begin{equation}
\mathbf x_n=
 [\Re G_{n1},\ldots,\Re G_{nK},
  \Im G_{n1},\ldots,\Im G_{nK}]^T.
\end{equation}
The full real vector is
$\mathbf x=[\mathbf x_1^T,\ldots,\mathbf x_N^T]^T\in\mathbb R^{2NK}$.
This ordering gives
$\|\widehat{\mathbf x}-\mathbf x\|_2^2
=\|\widehat{\mathbf G}-\mathbf G\|_F^2$.
We condition on the known $\mathbf S$, $\mathbf B$ and $\sigma^2$.
The noise convention is $\mathbb E|W_{np}|^2=\sigma^2$, so each real
component has variance $\sigma^2/2$.

Write $\mu_{np}=(\mathbf G\mathbf S+\mathbf B)_{np}$ and
$a_{np}=|\mu_{np}|$. For $z\geq0$, the amplitude density is
\begin{equation}
p(z\mid a)=\frac{2z}{\sigma^2}
 \exp\!\left(-\frac{z^2+a^2}{\sigma^2}\right)
 I_0\!\left(\frac{2za}{\sigma^2}\right).
\end{equation}
With $R(u)=I_1(u)/I_0(u)$, its amplitude score is
\begin{equation}
\frac{\partial\log p(z\mid a)}{\partial a}
=\frac{2}{\sigma^2}
 \left[zR\!\left(\frac{2za}{\sigma^2}\right)-a\right].
\end{equation}
Using $\mathbb E[zR(2za/\sigma^2)]=a$, the scalar information is
\begin{align}
\mathcal I_a(a)
&=\frac{4}{\sigma^4}
 \left[\mathbb E\!\left\{z^2R^2\!\left(\frac{2za}{\sigma^2}\right)\right\}
 -a^2\right],
\label{eq:revision-scalar-information}\\
\beta(a)&=\mathcal I_a(a)/4.
\end{align}

For $\mu_{np}\ne0$, put
\begin{equation}
\mathbf c_{np}=e^{-\jmath\arg\mu_{np}}\mathbf S_{:,p},
\qquad
\mathbf v_{np}=
\begin{bmatrix}\Re\mathbf c_{np}\\-\Im\mathbf c_{np}\end{bmatrix}.
\end{equation}
Here $\mathbf v_{np}=\partial a_{np}/\partial\mathbf x_n$.
Independence of the measurement noises gives
\begin{align}
\mathbf J_n&=\sum_{p=1}^{P}
  \mathcal I_a(a_{np})\mathbf v_{np}\mathbf v_{np}^T,\\
\mathbf J&=\operatorname{blkdiag}
 (\mathbf J_1,\ldots,\mathbf J_N).
\label{eq:revision-real-fisher}
\end{align}
Each nonzero measurement term has rank one in real coordinates. At
$\mu_{np}=0$, the contribution is defined by its continuous zero limit;
an arbitrary choice of phase supplies no information.

\subsection{Geometric tangent space}

For each user $k$ and path $\ell$, take the real parameters
$(\psi_{\ell,k},\Re\alpha_{\ell,k},\Im\alpha_{\ell,k})$.
Collect them in $\boldsymbol\eta$ and set
$\mathbf D=\partial\mathbf x/\partial\boldsymbol\eta$.
Writing $v_n(\psi)=e^{-\jmath(n-1)\psi}$, the complex derivatives for
one path are
\begin{align}
\frac{\partial g_k[n]}{\partial\psi_{\ell,k}}
 &=-\jmath(n-1)\alpha_{\ell,k}v_n(\psi_{\ell,k}),\\
\frac{\partial g_k[n]}{\partial\Re\alpha_{\ell,k}}
 &=v_n(\psi_{\ell,k}),\\
\frac{\partial g_k[n]}{\partial\Im\alpha_{\ell,k}}
 &=\jmath v_n(\psi_{\ell,k}).
\end{align}
Their real and imaginary parts fill the corresponding entries of
$\mathbf x_n$; entries belonging to other users are zero. Using spatial
frequency $\psi$ follows the channel generator directly and avoids a
coordinate singularity from $\psi=\pi\sin\theta$ at endfire.

Take a thin SVD $\mathbf D=\mathbf U_D\boldsymbol\Sigma_D\mathbf V_D^T$.
The implementation retains the columns of $\mathbf U_D$ whose singular
values exceed $10^{-9}$ times the largest singular value. These columns
form $\mathbf U$, an orthonormal basis for the numerically resolved range
of $\mathbf D$. This construction avoids inverting a redundant physical
parameterization when $\mathbf D$ is not full column rank.

For locally unbiased estimation along this tangent space,
\begin{equation}
\mathbf C_{\rm geom}
=\mathbf U(\mathbf U^T\mathbf J\mathbf U)^{-1}\mathbf U^T.
\end{equation}
The gains and spatial frequencies are varied jointly in $\mathbf D$;
neither is held fixed as known side information. This calculation includes
no unknown reference-field, calibration or noise-variance parameters.
Adding such nuisance parameters would require the corresponding enlarged
information matrix. The true path counts and parameter point specify the
local reference, rather than inputs available to HS-GS.

The displayed inverses require positive definite information on the
relevant space. A singular restricted matrix indicates an unresolved
local direction and is not replaced by a regularized finite bound.
The known reference can remove phase ambiguity, but its presence alone
does not guarantee identifiability for every pilot design. Coincident or
zero-gain paths can also invalidate a regular local parameterization.
The stated numerical rank tolerance must therefore accompany the results.

\subsection{Numerical evaluation and normalization}

For trial $i$, the two scalar reference errors are
\begin{align}
b_{{\rm unc},i}&=\sum_{n=1}^{N}\operatorname{tr}(\mathbf J_{n,i}^{-1}),\\
b_{{\rm geom},i}&=
 \operatorname{tr}\!\left[(\mathbf U_i^T\mathbf J_i\mathbf U_i)^{-1}\right].
\end{align}
The second identity uses $\mathbf U_i^T\mathbf U_i=\mathbf I$.
For either reference, the plotted value is
\begin{equation}
10\log_{10}
 \frac{\sum_i b_i}{\sum_i\|\mathbf G_i\|_F^2}.
\label{eq:revision-bound-normalization}
\end{equation}
The same channel realizations and energy denominators are used for the
estimator curves. This is a ratio of sums, not an average of trial-wise
dB values.

The scalar expectation is evaluated with scaled Bessel functions and
quadrature over $[\max(0,a-8\sigma),a+8\sigma]$, where
$\sigma=\sqrt{\sigma^2}$. The quadrature uses absolute tolerance
$10^{-14}$ and relative tolerance $10^{-12}$, splitting at $a$ when it
lies inside the integration window. At each main-figure SNR, 800 equally
spaced amplitude nodes cover $0.99a_{\min}$ to $1.01a_{\max}$ over that
SNR's worlds. Piecewise cubic Hermite interpolation supplies the scalar
weights at the measurement amplitudes. The implementation checks 200
sampled amplitudes against direct quadrature and requires maximum relative
error below $10^{-4}$. The tangent derivatives are checked by finite
differences. It also checks that the constrained trace is positive and
no larger than the unconstrained trace, and records the condition number
of the restricted information matrix. No diagonal ridge is added.
